\documentclass{article}

% packages
\usepackage{amsmath, amsthm, thmtools, amsfonts, amssymb, luacode, catchfile, tikzducks, hyperref, ifthen}
\ifcsname c@kobocompile\endcsname
	\usepackage[a5paper, total={1072pt, 1448pt}, margin=10pt, includeheadfoot]{geometry} % set page margins
\else
	\usepackage[a4paper, margin=50pt, includeheadfoot]{geometry}
\fi
\usepackage[shortlabels]{enumitem}
\usepackage[skip=3pt, indent=0pt]{parskip}

% language
\usepackage[bidi=basic, layout=tabular, provide=*]{babel}
\ifcsname c@english\endcsname
	\babelprovide[main, import]{english}
\else
	\babelprovide[main, import]{hebrew}
	\babelprovide{rl}
\fi
%\babelfont{rm}{Libertinus Serif}
\babelfont{rm}[Renderer=Harfbuzz]{Libertinus Serif}
\babelfont{sf}{Libertinus Sans}
\babelfont{tt}{Libertinus Mono}

% style
\AddToHook{cmd/section/before}{\clearpage}	% Add line break before section
\linespread{1.3}
\setcounter{secnumdepth}{0}		% Remove default number tags from sections, this won't do well with theorems
\AtBeginDocument{\setlength{\belowdisplayskip}{3pt}}
\AtBeginDocument{\setlength{\abovedisplayskip}{3pt}}
\graphicspath{ {../images/} }

% operators
\DeclareMathOperator\cis{cis}
\DeclareMathOperator\Sp{Sp}
\DeclareMathOperator\tr{tr}
\DeclareMathOperator\im{Im}
\DeclareMathOperator\re{Re}
\DeclareMathOperator\diag{diag}
\DeclareMathOperator*\lowlim{\underline{lim}}
\DeclareMathOperator*\uplim{\overline{lim}}
\DeclareMathOperator\rng{rng}
\DeclareMathOperator\Sym{Sym}
\DeclareMathOperator\Arg{Arg}
\DeclareMathOperator\Log{Log}
\DeclareMathOperator\dom{dom}
\DeclareMathOperator\supp{Supp}
\DeclareMathOperator\var{Var}
\DeclareMathOperator\cov{Cov}

% commands
%\renewcommand\qedsymbol{\textbf{מש''ל}}
%\renewcommand\qedsymbol{\fbox{\emoji{lizard}}}
\newcommand{\Aa}[0]{\mathcal{A}}
\newcommand{\Bb}[0]{\mathcal{B}}
\newcommand{\CC}[0]{\mathbb{C}}
\newcommand{\Cc}[0]{\mathcal{C}}
\newcommand{\EE}[0]{\mathbb{E}}
\newcommand{\FF}[0]{\mathbb{F}}
\newcommand{\Ff}[0]{\mathcal{F}}
\newcommand{\Ii}[0]{\mathcal{I}}
\newcommand{\Gg}[0]{\mathcal{G}}
\newcommand{\Ll}[0]{\mathcal{L}}
\newcommand{\Mm}[0]{\mathcal{M}}
\newcommand{\NN}[0]{\mathbb{N}}
\newcommand{\Nn}[0]{\mathcal{N}}
\newcommand{\PP}[0]{\mathbb{P}}
\newcommand{\Pp}[0]{\mathcal{P}}
\newcommand{\QQ}[0]{\mathbb{Q}}
\newcommand{\RR}[0]{\mathbb{R}}
\newcommand{\Rr}[0]{\mathcal{R}}
\newcommand{\Ss}[0]{\mathcal{S}}
\newcommand{\TT}[0]{\mathbb{T}}
\newcommand{\Uu}[0]{\mathcal{U}}
\newcommand{\Vv}[0]{\mathcal{V}}
\newcommand{\Ww}[0]{\mathcal{W}}
\newcommand{\ZZ}[0]{\mathbb{Z}}
\newcommand{\acts}[0]{\circlearrowright}
\newcommand{\explain}[2] {
	\begin{flalign*}
		 && \text{#2} && \text{#1}
	\end{flalign*}
}
\newcommand{\maketitleprint}[0]{ \begin{center}
	%\begin{tikzpicture}[scale=3]
	%	\duck[graduate=gray!20!black, tassel=red!70!black]
	%\end{tikzpicture}	
	\includegraphics[width=6cm]{cover}
\end{center}
}

% theorem commands
\newtheoremstyle{c_remark}
	{}	% Space above
	{}	% Space below
	{}% Body font
	{}	% Indent amount
	{\bfseries}	% Theorem head font
	{}	% Punctuation after theorem head
	{.5em}	% Space after theorem head
	{\thmname{#1}\thmnumber{ #2}\thmnote{ \normalfont{\text{(#3)}}}}	% head content
\newtheoremstyle{c_definition}
	{3pt}	% Space above
	{3pt}	% Space below
	{}% Body font
	{}	% Indent amount
	{\bfseries}	% Theorem head font
	{}	% Punctuation after theorem head
	{.5em}	% Space after theorem head
	{\thmname{#1}\thmnumber{ #2}\thmnote{ \normalfont{\text{(#3)}}}}	% head content
\newtheoremstyle{c_plain}
	{3pt}	% Space above
	{3pt}	% Space below
	{\itshape}% Body font
	{}	% Indent amount
	{\bfseries}	% Theorem head font
	{}	% Punctuation after theorem head
	{.5em}	% Space after theorem head
	{\thmname{#1}\thmnumber{ #2}\thmnote{ \text{(#3)}}}	% head content

\ifcsname c@english\endcsname
	\theoremstyle{plain}
	\newtheorem{theorem}{Theorem}[section]
	\newtheorem{lemma}[theorem]{Lemma}
	\newtheorem{proposition}[theorem]{Proposition}
	\newtheorem*{proposition*}{Proposition}
	%\newtheorem{corollary}[theorem]{אין חלופה עברית}

	\theoremstyle{definition}
	\newtheorem{definition}[theorem]{Definition}
	\newtheorem*{definition*}{Definition}
	\newtheorem{example}{Example}[section]
	\newtheorem{exercise}{Exercise}[section]

	\theoremstyle{remark}
	\newtheorem*{remark}{Remark}
	\newtheorem*{solution}{Solution}
	\newtheorem{conclusion}[theorem]{Conclusion}
	\newtheorem{notation}[theorem]{Notation}
\else
	\theoremstyle{c_plain}
	\newtheorem{theorem}{משפט}[section]
	\newtheorem{lemma}[theorem]{למה}
	\newtheorem{proposition}[theorem]{טענה}
	\newtheorem*{proposition*}{טענה}
	%\newtheorem{corollary}[theorem]{אין חלופה עברית}

	\theoremstyle{c_definition}
	\newtheorem{definition}[theorem]{הגדרה}
	\newtheorem*{definition*}{הגדרה}
	\newtheorem{example}{דוגמה}[section]
	\newtheorem{exercise}{תרגיל}[section]

	\theoremstyle{c_remark}
	\newtheorem*{remark}{הערה}
	\newtheorem*{solution}{פתרון}
	\newtheorem{conclusion}[theorem]{מסקנה}
	\newtheorem{notation}[theorem]{סימון}
\fi

% Questions related commands
\newcounter{question}
\setcounter{question}{1}
\newcounter{sub_question}
\setcounter{sub_question}{1}

\ifcsname c@english\endcsname
	\newcommand{\question}[1][0]{
		\ifthenelse{#1 = 0}{}{\setcounter{question}{#1}}
		\section{Question \arabic{question}}
		\addtocounter{question}{1}
		\setcounter{sub_question}{1}
	}

	\newcommand{\subquestion}[1][0]{
		\ifthenelse{#1 = 0}{}{\setcounter{sub_question}{#1}}
		\subsection{Part \alph{sub_question}}
		\addtocounter{sub_question}{1}
	}
\else
	\newcommand{\question}[1][0]{
		\ifthenelse{#1 = 0}{}{\setcounter{question}{#1}}
		\section{שאלה \arabic{question}}
		\addtocounter{question}{1}
		\setcounter{sub_question}{1}
	}

	\newcommand{\subquestion}[1][0]{
		\ifthenelse{#1 = 0}{}{\setcounter{sub_question}{#1}}
		\subsection{סעיף \localecounter{letters.gershayim}{sub_question}}
		\addtocounter{sub_question}{1}
	}
\fi

% import lua and start of document
\directlua{common = require ('../common')}

\GetEnv{AUTHOR}

% headers
\author{\AUTHOR}
\date\today


\title{פתרון ממ''ן 12 – אלגברה לינארית 2 (20229)}

\begin{document}
\maketitle

\section{שאלה 1}
\subsection{סעיף א'}
נגדיר
\[
	A_1 = \begin{pmatrix}
		0 & i \\
		-i & 0
	\end{pmatrix},
	A_2 = \begin{pmatrix}
		i & 2i \\
		0 & 1
	\end{pmatrix},
	A_3 = \begin{pmatrix}
		1 & i \\
		1 & 2 + i
	\end{pmatrix}
\]
נבדוק אם כל אחת מן המטריצות נורמליות ואם כן נמצא מטריצה אוניטרית המלכסנת אותן. \\*
נבדוק אם $A_1$ נורמלית:
\[
	A_1^* = \begin{pmatrix}
		0 & i \\
		-i & 0
	\end{pmatrix} = A_1
\]
לכן בהכרח $A_1 A_1^* = A_1^* A_1$ וכן $A_1$ נורמלית.
נחשב את הפולינום האופייני של $A_1$:
\[
	p(t) = \begin{vmatrix}
		t & -i \\
		i & t
	\end{vmatrix}
	= t^2 - (-i)i = t^2 - 1 = (t - 1)(t + 1)
\]
אז למטריצה שני ערכים עצמיים $-1, 1$.
נמצא את $V_{-1}$ על־ידי פתרון המערכת:
\[
	A_1 u = -u \rightarrow (A_1 + I) u = 0 \rightarrow \begin{pmatrix}
		1 & i \\
		-i & 1
	\end{pmatrix}
	\xrightarrow{R_2 \to R_2 + iR_1}
	\begin{pmatrix}
		1 & i \\
		0 & 0
	\end{pmatrix}
	\rightarrow
	V_{-1} = \Sp\{ (-i, 1) \}
\]
לכן בבסיס האורתוגונלי של $V_{-1}$ ישנו רק הווקטור $\begin{pmatrix} \frac{-i}{\sqrt{2}} \\ \frac{1}{\sqrt{2}} \end{pmatrix}$. \\*
נמצא את המרחב $V_1$ באופן דומה:
\[
	(A_1 - I) = 0
	\rightarrow
	\begin{pmatrix}
		-1 & i \\
		-i & -1
	\end{pmatrix}
	\xrightarrow[R_2 \to i R_2, R_2 \to R_2 - R_1]{R_1 \to -R_1}
	\begin{pmatrix}
		1 & -i \\
		0 & 0
	\end{pmatrix}
	\rightarrow
	V_1 = \Sp\{ (i, 1) \}
\]
בבסיס האורתוגונלי של $V_{-1}$ ישנו הווקטור $\begin{pmatrix} \frac{i}{\sqrt{2}} \\ \frac{1}{\sqrt{2}} \end{pmatrix}$. \\*
בשל כך המטריצה $P$ המוגדרת להלן מלכסנת אוניטרית את $A_1$ לפי משפט 3.2.1:
\[
	P = \begin{pmatrix}
		\frac{-i}{\sqrt{2}} & \frac{i}{\sqrt{2}} \\
		\frac{1}{\sqrt{2}} & \frac{1}{\sqrt{2}}
	\end{pmatrix}
\]
נחשב את $A_2^*$:
\[
	A_2^* = \begin{pmatrix}
		-i & 0 \\
		-2i & 1
	\end{pmatrix}
\]
נבדוק אם $A_2$ נורמלית:
\begin{align*}
	A_2 A_2^* = 
	\begin{pmatrix}
		i & 2i \\
		0 & 1
	\end{pmatrix} \begin{pmatrix}
		-i & 0 \\
		-2i & 1
	\end{pmatrix}
	= \begin{pmatrix}
		3 & 2i \\
		-2i & 1
	\end{pmatrix} \\
	A_2^* A_2 = 
	\begin{pmatrix}
		-i & 0 \\
		-2i & 1
	\end{pmatrix}
	\begin{pmatrix}
		i & 2i \\
		0 & 1
	\end{pmatrix}
	= \begin{pmatrix}
		1 & 2 \\
		2 & 5
	\end{pmatrix} \\
\end{align*}
אנו רואים כי $A_2 A_2^* \ne A_2^* A_2$ ולכן $A_2$ איננה נורמלית. \\*
נבדוק אם $A_3$ נורמלית, תחילה נבחין כי
\[
	A_3^* = \begin{pmatrix}
		1 & 1 \\
		-i & 2 - i
	\end{pmatrix}
\]
נבדוק אם היא נורמלית:
\begin{align*}
	A_3 A_3^* = 
	\begin{pmatrix}
		1 & i \\
		1 & 2 + i
	\end{pmatrix}
	\begin{pmatrix}
		1 & 1 \\
		-i & 2 - i
	\end{pmatrix}
	= \begin{pmatrix}
		2 & 2 + 2i \\
		2 - 2i & 6
	\end{pmatrix} \\
	A_3^* A_3 = 
	\begin{pmatrix}
		1 & 1 \\
		-i & 2 - i
	\end{pmatrix}
	\begin{pmatrix}
		1 & i \\
		1 & 2 + i
	\end{pmatrix}
	= \begin{pmatrix}
		2 & 2 + 2i \\
		2 - 2i & 6
	\end{pmatrix}
\end{align*}
אנו רואים כי המטריצה $A_3$ אכן נורמלית. \\*
נחשב את הפולינום האופייני של $A_3$:
\[
	p(t) = \begin{vmatrix}
		t - 1 & -i \\
		-1 & t - 2 - i
	\end{vmatrix}
	= (t - 1)(t - 2 - i) + i
	= t^2 + (-3 - i)t + 2
	\rightarrow
	t = \frac{3 + i \pm \sqrt{8 + 6i - 8}}{2}
	= \frac{3 + i \pm (\sqrt{3} + \sqrt{3}i)}{2}
\]
נשתמש בחישוב שמבוצע בתשובה 3.2.2 עבור המטריצה ולכן המטריצה האוניטרית היא
\[
	\begin{pmatrix}
		\frac{1}{\sqrt{3 + \sqrt{3}}} & \frac{1}{\sqrt{3 - \sqrt{3}}} \\
		\frac{1 + \sqrt{3}}{2\sqrt{3 + \sqrt{3}}}(1 - i) & \frac{1 - \sqrt{3}}{2\sqrt{3 - \sqrt{3}}}(1 - i)
	\end{pmatrix}
\]

\subsection{סעיף ב'}
נמצא אילו מבין המטריצות המוגדרות להלן חיוביות. \\*
נגדיר
\[
	C_1 = \begin{pmatrix}
		1 & 1 \\
		1 & 1
	\end{pmatrix}
\]
המטריצה ממשית וסימטרית, ולכן גם נורמלית. נחשב את ערכיה העצמיים:
\[
	p(t) = {(t - 1)}^2 - 1 = t^2 - 2t = t(t - 2)
\]
לכן ערכיה העצמיים הם $0, 2$ והיא חיובית אך לא חיובית לחלוטין. \\*
נגדיר
\[
	C_2 = \begin{pmatrix}
		0 & i \\
		-i & 0
	\end{pmatrix}
\]
קל לראות כי המטריצה צמודה לעצמה, נחשב את ערכיה העצמיים:
\[
	p(t) = t^2 - 1 = (t - 1)(t + 1)
\]
למטריצה ערך עצמי שלילי ולכן היא לא חיובית כלל לפי משפט 3.2.2.
נגדיר
\[
	C_3 = \begin{pmatrix}
		0 & 1 \\
		-1 & 0
	\end{pmatrix}
\]
מטריצה זו לא צמודה לעצמה, ולכן לא יכולה להיות מטריצה חיובית כלל. \\*
נגדיר
\[
	C_4 = \begin{pmatrix}
		1 & 1 \\
		0 & 1
	\end{pmatrix}
\]
גם מטריצה זו לא סימטרית ולכן לא עומדת בהגדרה 1.2.5 כלל. \\*
נגדיר
\[
	C_5 = \begin{pmatrix}
		2 & 1 \\
		1 & 2
	\end{pmatrix}
\]
מטריצה זו סימטרית, נחשב את ערכיה העצמיים:
\[
	p(t) = {(t - 2)}^2 - 1 = t^2 - 4t + 4 - 1 = t^2 - 4t + 3 = (t - 3)(t - 1)
\]
לכן לפי משפט 3.2.2 המטריצה חיובית לחלוטין. \\*
נגדיר
\[
	C_6 = \begin{pmatrix}
		1 & 2 \\
		3 & 1
	\end{pmatrix}
\]
גם מטריצה זו לא סימטרית ולכן לא חיובית.

\section{שאלה 2}
יהי $V$ מרחב מכפלה פנימית מממד סופי ותהי $T: V \to V$ העתקה לינארית נורמלית.
\subsection{i}
נוכיח כי $\ker T = \ker T^*$. \\*
יהי $u \in \ker T$, אז מתקיים $T u = 0$ ובשל כך גם $\lVert u \rVert = 0$. \\*
על־פי הגדרת הנורמה מתקיים גם $(T u, T u) = 0$. אז
\begin{align*}
	\lVert T u \rVert & = (T u, T u) \\
					  & = (u, T^* T u) && \text{מהגדרת הצמוד} \\
					  & = (u, T T^* u) && \text{ידוע כי $T$ נורמלית} \\
					  & = (T^* u, T^* u) && \text{שוב על־פי צמוד} \\
					  & = \lVert T^* u \rVert = 0 && \text{הגדרת הנורמה} \\
					  & \rightarrow T^* u = 0 \rightarrow u \in \ker T^* && \text{נורמה חיובית לחלוטין}
\end{align*}
בשל הגדרת המכפלה הפנימית נובע כי $T^* u = 0$. \\*
נשים לב כי בשל סימטריות המטריצות הנורמליות נוכל להוכיח באופן דומה גם כי אם $T^* u = 0$ אז $T u = 0$,
ולכן מתקיים $\ker T = \ker T^*$.

\subsection{ii}
נוכיח כי $\im T = {(\ker T)}^\perp$. \\*
על־פי משפט 3.2.1 ההעתקה $T$ היא לכסינה אוניטרית בפרט ולכן לכסינה בכלל.
מסיבה זו כל וקטור ב־$V$ הוא וקטור עצמי לאיזשהו ערך עצמי ב־$T$. \\*
יהיו $u \in \im T, v \in \ker T$. נשים לב כי $v$ הוא וקטור עצמי של $0$, ואילו $u$ וקטור עצמי לערך עצמי $\lambda \ne 0$. \\*
בשל היותם ערכים עצמיים שונים, לפי משפט 3.2.6 הוקטורים אורתוגונליים זה לזה. \\*
נכליל את הטענה הזו ונראה שלכל $u$ כזה התנאי מתקיים לכל $v$, לכן $u \in {(\ker T)}^\perp$. \\*
באופן דומה נוכל להכליל את הטענה על התמונה, ולכן $\im T = {(\ker T)}^\perp$.

\subsection{iii}
נוכיח כי $\im T = \im T^*$. \\*
על־פי סעיף ii מתקיים
\[
	\im T^* = {(\ker T^*)}^\perp
\]
וידוע כי $\ker T = \ker T^*$ על־פי סעיף i, לכן
\[
	\im T^* = {(\ker T)}^\perp = \im T
\]

\section{שאלה 3}
יהי $V$ מרחב מכפלה עצמית מממד סופי ו־$T: V \to V$ העתקה לינארית המקיימת
\[
	T^2 = \frac{1}{2}(T + T^*)
\]
נוכיח כי $T$ נורמלית ומתקיים $T^2 = T$.
\[
	T T^* = T (2T^2 - T) = T (2T - I) T = (2T^2 - T) T = T^* T
\]
אנו רואים כי $T$ נורמלית. \\*
נוכיח גם כי $T^2 = T$. \\*
יהי $\lambda$ ערך עצמי של $T$ ו־$u$ וקטור עצמי שלו. \\*
לפי למה 3.2.5 $\overline{\lambda}$ ערך עצמי של $T^*$ ו־$u$ וקטור עצמי שלו,
דהינו $T u = \lambda u, T^* u = \overline{\lambda} u$. \\*
עוד אנו יודעים כי $T^2 u = T \lambda u = \lambda^2 u$. מתקיים
\[
	T^2 u = \lambda^2 u = \frac{1}{2}(\lambda u + \overline{\lambda} u) = \frac{1}{2}(T u + T^* u)
\]
אז נוכל להניח גם כי
\[
	2 \lambda^2 = \lambda + \overline{\lambda}
\]
נניח כי המרחב מוגדר מעל $\CC$, כך שכל מרחב מכפלה פנימית הלכה למעשה מוכל בו.
נגדיר $\lambda = a + bi$, אז
\[
	2a^2 + 4abi - 2b^2 = a + bi + a - bi
\]
דהינו מתקיים
\[
	\begin{cases}
		2ab = 0 \\
		2a^2 - 2b^2 = 2a
	\end{cases}
\]
אילו $b \ne 0$, אז נובע מהמשוואה הראשונה כי $a = 0$ ומהשמוואה השנייה כי $b = 0$ בסתירה לטענה, אז מהמשוואה הראשונה אנו למדים כי $a \ne 0$.
לכן גם $b = 0$ ומהמשוואה השנייה נובע ש־$a(a - 1) = 0$, ידוע כי $a \ne 0$ ולכן $a = 1$.
מצאנו כי כלל הערכים העצמיים של $T, T^2, T^*$ הם $1$ בלבד, לכן לכל $v \in V$ מתקיים $T v = T^2 v = v$ ובהתאם $T^2 = T$.

\section{שאלה 4}
תהי $H$ מטריצה סימטרית ממשית מסדר $n \times n$ ויהי $\lambda$ הערך העצמי המקסימלי של $H$. \\*
נוכיח כי לכל $v \in \RR^n$ כאשר $\lVert v \rVert = 1$ מתקיים $v^t H v \le \lambda$.
\begin{proof}
	מהגדרה 2.3.4 נובע כי $H$ אוניטרית בכלל ולכן נורמלית בפרט ולכן לפי הגדרה 3.2.1 $H$ לכסינה אוניטרית. \\*
	נגדיר $B = (b_1, \hdots, b_n)$ בסיס מלכסן אורתונורמלי של $H$ כך שמתקיים לכל $u \in V$ עבורו ${[u]}_B = (\alpha_1, \hdots, \alpha_n)$:
	\[
		H u = \sum_{i = 1}^n \lambda_i \alpha_i b_i \tag{1}
	\]
	יהי $v \in V$ וקטור כך ש־$\lVert v \rVert = 1$. ממשפט 2.3.2 נובע כי
	\begin{align*}
		{\lVert v \rVert}^2 = 1^2
		& = {\lVert H v \rVert}^2 = (H v, H v) \\
		& = (\sum_{i = 1}^n \lambda_i \alpha_i b_i, \sum_{j = 1}^n \lambda_j \alpha_j b_j) && \text{על־פי (1)} \\
		& = \sum_{i = 1}^n \sum_{j = 1}^n \lambda_j \alpha_j \lambda_i \alpha_i (b_i, b_j) && \text{המכפלה הפנימית הסטנדרטית} \\
		& = \sum_{i = 1}^n \lambda_i^2 \alpha_i^2 && \text{$B$ בסיס אורתונורמלי} \\
		1 & = \sum_{i = 1}^n \alpha_i^2 && \text{טענה 2.4.3} \tag{2}
	\end{align*}
	נגדיר $\lambda = \max\{ \lambda_i \mid 1 \le i \le n \}$ ונראה כי
	\begin{align*}
		v^t H v & = (v, H v) \\
				& = (\sum_{i = 1}^n \alpha_i b_i, \sum_{j = 1}^n \lambda_j \alpha_j b_j) \\
		= \sum_{i = 1}^n \lambda_i \alpha_i^2 & \le \sum_{i = 1}^n \lambda \alpha_i^2 && \lambda \ge \lambda_i \\
		v^t H v & \le \lambda \sum_{i = 1}^n \alpha_i^2 \\
		v^t H v & \le \lambda && \text{על־פי (2)}
	\end{align*}
\end{proof}

\section{שאלה 5}
נוכיח כי המטריצה $A$ המוגדרת להלן היא נורמלית.
\[
	A = \begin{pmatrix}
		2 - i & -1 & 0 \\
		-1 & 1 - i & 1 \\
		0 & 1 & 2 - i
	\end{pmatrix}
\]
לפי הגדרת המשלים
\[
	A^* = \begin{pmatrix}
		2 + i & -1 & 0 \\
		-1 & 1 + i & 1 \\
		0 & 1 & 2 + i
	\end{pmatrix}
\]
נחשב
\begin{align*}
	A A^* = 
	\begin{pmatrix}
		2 - i & -1 & 0 \\
		-1 & 1 - i & 1 \\
		0 & 1 & 2 - i
	\end{pmatrix}
	\begin{pmatrix}
		2 + i & -1 & 0 \\
		-1 & 1 + i & 1 \\
		0 & 1 & 2 + i
	\end{pmatrix}
	= \begin{pmatrix}
		6 & -3 & -1 \\
		-3 & 4 & 3 \\
		-1 & 3 & 6
	\end{pmatrix} \\
	A^* A = 
	\begin{pmatrix}
		2 + i & -1 & 0 \\
		-1 & 1 + i & 1 \\
		0 & 1 & 2 + i
	\end{pmatrix}
	\begin{pmatrix}
		2 - i & -1 & 0 \\
		-1 & 1 - i & 1 \\
		0 & 1 & 2 - i
	\end{pmatrix}
	= \begin{pmatrix}
		6 & -3 & -1 \\
		-3 & 4 & 3 \\
		-1 & 3 & 6
	\end{pmatrix}
\end{align*}
על־פי החיושב $A A^* = A^* A$ ולכן המטריצה נורמלית. \\*
נמצא את הפירוק
\[
	A = \sum_{i} \lambda_i P_i
\]
כאשר $P_i$ הן מטריצות של ההטלות האורתוגונליות שמופיעות בפירוק הספקטראלי של $T_A$. \\*
נחשב את הפולינום האופייני של $A$:
\begin{align*}
	p(t) 
	& = 
	\begin{vmatrix}
		t - 2 + i & 1 & 0 \\
		1 & t - 1 + i & -1 \\
		0 & -1 & t - 2 + i
	\end{vmatrix}
	\xrightarrow{R_3 \to R_3 + R_1} 
	\begin{vmatrix}
		t - 2 + i & 1 & 0 \\
		1 & t - 1 + i & -1 \\
		t - 2 + i & 0 & t - 2 + i
	\end{vmatrix} \\
	& = 
	(t - 2 + i)
	\begin{vmatrix}
		t - 2 + i & 1 & 0 \\
		1 & t - 1 + i & -1 \\
		1 & 0 & 1
	\end{vmatrix}
	\xrightarrow{R_2 \to R_2 + R_3}
	(t - 2 + i)
	\begin{vmatrix}
		t - 2 + i & 1 & 0 \\
		2 & t - 1 + i & 0 \\
		1 & 0 & 1
	\end{vmatrix} \\
	& =
	(t - 2 + i)
	\begin{vmatrix}
		t - 2 + i & 1 \\
		2 & t - 1 + i \\
	\end{vmatrix}
	= (t - 2 + i) ((t - 2 + i) (t - 1 + i) - 2) \\
	& = (t - 2 + i)(t^2 + (-3 + 2i) t + (-1 - 3i))
	= (t - 2 + i)(t^2 + (-1 + 2i) t + 1 - 3i) \\
	& = (t - 2 + i)(t + i)(t - 3 + i)
\end{align*}
אז הערכים העצמיים של $T_A$ הם $2 - i, -i, 3 - i$. נחשב את המרחב העצמי שלהם על־ידי פתרון מערכת המשוואות $(A - \lambda I)u = 0$:
\begin{align*}
	& (A - (2 - i) I) u = 0
	\rightarrow
	\begin{pmatrix}
		0 & -1 & 0 \\
		-1 & -1 & 1 \\
		0 & 1 & 0
	\end{pmatrix}
	\xrightarrow[R_3 \to R_3 + R_1]{R_2 \to -R_2 + R_1}
	\begin{pmatrix}
		0 & -1 & 0 \\
		1 & 0 & -1 \\
		0 & 0 & 0
	\end{pmatrix} \\
	& \rightarrow
	V_{2 -i} = \Sp \{ (1, 0, 1) \}, v_1 = (1, 0, 1) \\
	& (A + i I) u = 0
	\rightarrow
	\begin{pmatrix}
		2 & -1 & 0 \\
		-1 & 1 & 1 \\
		0 & 1 & 2
	\end{pmatrix}
	\xrightarrow[R_1 \to -R_1]{R_2 \leftrightarrow R_1}
	\begin{pmatrix}
		1 & -1 & -1 \\
		2 & -1 & 0 \\
		0 & 1 & 2
	\end{pmatrix}
	\xrightarrow[R_3 \to R_3 - R_2 + 2R_1]{R_2 \to R_2 - 2R_1}
	\begin{pmatrix}
		1 & -1 & -1 \\
		0 & 1 & 2 \\
		0 & 0 & 0
	\end{pmatrix} \\
	& \rightarrow
	V_{-i} = \Sp\{ {(1, 2, -1)}^t \}, v_2 = (1, 2, -1) \\
	& (A - (3 - i) I) u = 0
	\rightarrow
	\begin{pmatrix}
		-1 & -1 & 0 \\
		-1 & -2 & 1 \\
		0 & 1 & -1
	\end{pmatrix}
	\xrightarrow[R_2 \to R_2 + R_1]{R_1 \to - R_1}
	\begin{pmatrix}
		1 & 1 & 0 \\
		0 & -1 & 1 \\
		0 & 1 & -1
	\end{pmatrix}
	\rightarrow
	\begin{pmatrix}
		1 & 1 & 0 \\
		0 & 1 & -1 \\
		0 & 0 & 0
	\end{pmatrix} \\
	& \rightarrow V_{3 - i} = \Sp\{ {(1, -1, -1)}^t \}, v_3 = (1, -1, -1)
\end{align*}
נגדיר $B = (v_1, v_2, v_3)$ בסיס אורתוגונלי ו־$E$ הבסיס הסטנדרטי, אז מטריצת המעבר מ־$E$ ל־$B$ היא
\[
	M = \begin{pmatrix}
		1 & 1 & 1 \\
		0 & 2 & -1 \\
		1 & -1 & -1
	\end{pmatrix}
\]
לפי חישוב גם
\[
	M^{-1} = \begin{pmatrix}
		\frac{1}{2} & 0 & \frac{1}{2} \\
		-\frac{1}{2} & 1 & \frac{1}{2} \\
		1 & -1 & -1
	\end{pmatrix}
\]
נגדיר $P_i$ ההטלה האורתוגונלית על האיבר ה־$i$, אז לפי הקורס הקודם
\[
	{[P_i]}_E = M {[P_i]}_B M^{-1}
\]
\begin{align*}
	& {[P_1]}_E = 
	\begin{pmatrix}
		1 & 1 & 1 \\
		0 & 2 & -1 \\
		1 & -1 & -1
	\end{pmatrix}
	\begin{pmatrix}
		1 & 0 & 0 \\ 
		0 & 0 & 0 \\ 
		0 & 0 & 0
	\end{pmatrix}
	\begin{pmatrix}
		\frac{1}{2} & 0 & \frac{1}{2} \\
		-\frac{1}{2} & 1 & \frac{1}{2} \\
		1 & -1 & -1
	\end{pmatrix}
	= \begin{pmatrix}
		\frac{1}{2} & 0 & 0 \\
		0 & 0 & 0 \\
		1 & 0 & 0
	\end{pmatrix} \\
	& {[P_2]}_E = 
	\begin{pmatrix}
		1 & 1 & 1 \\
		0 & 2 & -1 \\
		1 & -1 & -1
	\end{pmatrix}
	\begin{pmatrix}
		0 & 0 & 0 \\ 
		0 & 1 & 0 \\ 
		0 & 0 & 0
	\end{pmatrix}
	\begin{pmatrix}
		\frac{1}{2} & 0 & \frac{1}{2} \\
		-\frac{1}{2} & 1 & \frac{1}{2} \\
		1 & -1 & -1
	\end{pmatrix}
	= \begin{pmatrix}
		0 & 0 & 0 \\
		0 & 2 & 0 \\
		0 & 1 & 0
	\end{pmatrix} \\
	& {[P_3]}_E = 
	\begin{pmatrix}
		1 & 1 & 1 \\
		0 & 2 & -1 \\
		1 & -1 & -1
	\end{pmatrix}
	\begin{pmatrix}
		0 & 0 & 0 \\ 
		0 & 0 & 0 \\ 
		0 & 0 & 1
	\end{pmatrix}
	\begin{pmatrix}
		\frac{1}{2} & 0 & \frac{1}{2} \\
		-\frac{1}{2} & 1 & \frac{1}{2} \\
		1 & -1 & -1
	\end{pmatrix}
	= \begin{pmatrix}
		0 & 0 & \frac{1}{2} \\
		0 & 0 & \frac{1}{2} \\
		0 & 0 & 1
	\end{pmatrix} \\
\end{align*}
ולכן $A = (2 - i) {[P_1]}_E - i {[P_2]}_E + (3 - i) {[P_3]}_E$

\end{document}
